\fancyhead[L]{\small SEYFFER ON GEOGRAPHICAL LONGITUDE.}
\hypertarget{native-538Seyffer}{}\label{native-538Seyffer}
\begin{center}\small Göttingische gelehrte Anzeigen. No. 78. pp. 773..775. May 18, 1811.\\[3pt]\rule{.72\linewidth}{.3pt}\end{center}

\emph{On the Geographical Longitude of the Royal Observatory at Munich, First Determined from Thirty-Seven Solar Eclipses Observed and Worked Out by Calculation, by} \textsc{Carl Felix Seyffer}. 1810. 104 pages in quarto. Munich.

The title would lead one to suppose that the longitude was derived from 37 solar eclipses. In this paper, printed separately from the Memoirs of the Royal Bavarian Academy of Sciences, Mr. \textsc{Seyffer} instead derives the longitude of Munich by comparing a Munich observation of a single solar eclipse (June 16, 1806) with 36 observations made elsewhere. The places compared are \emph{Rome} (two observations), \emph{Padua, Milan, Madrid, Aranjues, Pampelona, Rinderhook, Fort Orange} (both in the United States of North America), \emph{Amsterdam, Utrecht,} \SourceLeaf{539}\emph{Zürich, Ochsenhausen, Leipzig, Breslau, Ofen, Cracau, Erlau} (in Hungary), \emph{Schweidnitz, Hamburg, Luck} (in Poland), \emph{Bourg en Bresse}, the island of \emph{Leon} (near Cadiz), \emph{Montauban, Toulouse, Paris, Prag, Lilienthal, Reikevich} (in Iceland), \emph{Göttingen, Naples, Brünn. Berlin, Regensburg, Cremsmünster.} To determine the longitude of Munich alone, it would of course have been unnecessary to take so many places into account. The work is valuable, however, as a contribution to determining the longitudes of those places themselves, and Mr. \textsc{Seyffer} has had all his calculations printed in full detail.

After using the complete observations to examine the corrections to the Moon's latitude, the horizontal parallax and the sum of the semidiameters, Mr. \textsc{Seyffer} finds the longitude of Munich to be $37^{\mathrm m}5^{\mathrm s}6$ east of Paris, as the mean of 15 comparisons. In the latest volume of the Connoissance des tems, Munich is listed among the places whose longitudes Mr. \textsc{Burckhardt} has examined afresh, and its longitude is given as $37^{\mathrm m}0^{\mathrm s}$.
